\section{考察}

本節では，主にヒアリングから得られた結果より，音楽嗜好を通じた自己開示の効果や，自己開示レベルの設定，推薦システムについて考察を行う．

\subsection{自己開示に関する考察}

本項では，音楽嗜好の共有が自己開示に与える影響について自己開示レベルの推移や，被験者の意見から考察を行う．

\subsubsection{自然な自己開示の傾向}
実験結果から，自由に選曲を行ったとき（推薦なし群）の自己開示の特徴的な傾向が明らかになった．表\ref{群別の自己開示レベルごとの曲数と人数}より，レベル4の自己開示が限定的であった一方で，レベル3の自己開示が頻繁に観察され，OとM以外の被験者全員がEQ5（聴いた曲について楽しく話ができた），EQ6（思いや経験が理解できた）で高スコア（4または5）を示した．この結果は，自己開示レベル3の曲が受容者にとって受け入れやすい曲であり，効果的な会話のきっかけとなったことを示唆している．

実験設計において曲に関する思い入れや経験の共有を主目的としたことで，自然と深い自己開示が促進されたと考えられる．レベル4の出現頻度が低かった主な要因として，曲の分類時に自己開示レベルの基準が不適切であった可能性が挙げられる．事後のヒアリングでは，ほとんどの被験者が強い思い入れのある曲を持っていたと報告しており，高いレベルの自己開示に該当するような曲があったにも関わらず，レベルの分類が不適切であった可能性がある．特に，「思い入れ」と「自己開示レベル」の関係性について，より精緻な基準設定が必要である．


\subsubsection{音楽嗜好の共有による自己開示}
楽曲に紐づいた個人的な経験や思い出が，自然な形で自己開示を促進することが確認された．例えば，高校時代の思い出や人生における重要な出来事，さらには恋愛体験などの個人的な経験が，楽曲を介して自然に共有されていた．これは音楽が「思い出の媒体」として機能し，通常の会話では語られにくい個人的な経験を引き出す触媒となることを示唆している．

音楽嗜好が自己開示を促す理由の1つに，長期的な友人関係においても，音楽を通じた新たな相互理解が促進されることが確認された．高校時代からの友人で頻繁にフェスに参加する被験者I,Jのペアでさえ，「こんな音楽も聴くんだってことを知れた」「自分の好きなものをたくさん伝えることができてよかった」と報告している．このことは，既知の関係性においても，音楽体験の共有が新たな一面の発見につながることを示唆している．
これらの結果から，音楽嗜好の共有は深い自己開示を促進する効果的な手段であることが実証された．音楽という媒体を介することで，より自然な形での自己開示が可能になると考えられる．


\begin{comment}
\subsubsection{楽曲への思い入れの影響}
楽曲に対する思い入れの強さが，自己開示の促進に強く結びついていることが明らかになった．思い入れの強い楽曲について語る際，より具体的かつ詳細な自己開示を行う傾向があることが分かった．
% 例えば，思い入れの強い楽曲では「年末のライブでこの曲を聴いた時の体験」や「この曲のここの歌い方が好きで」といった具体的な文脈を伴った自己開示が行われた．対照的に，思い入れの弱い楽曲では「知っている/知らない」や「他にどんな曲を出しているの？」といった表層的な会話に留まる傾向が見られた．
この結果は，楽曲への思い入れが強いほど，その曲に紐づいた個人的な体験や感情が豊富に存在し，それらの経験や感情が自然な会話の流れの中で語られやすくなることを示唆している．

音楽への関心度による影響も特徴的であった．音楽に対して深い関心を持つペアでは，楽曲の技術的な側面から個人的な体験まで，幅広く会話が展開された．一方で，音楽への関与度が比較的低いペアであっても，思い出という観点から会話が展開され，自己開示が自然に促進されることが確認された．
\end{comment}


\subsection{セッションについての考察}

本項では，音楽共有セッション自体の有効性について，被験者の意見や実験後のアンケートから分析を行う．特に，従来のカラオケとは異なる音楽共有セッションがもたらした効果について考察する．

\subsubsection{音楽を通じた会話をする楽しさ}
音楽共有セッションに対する被験者の反応は肯定的であった．19名の被験者が「楽しかった」と回答し，そのうち12名は「時間を忘れるほど楽しかった」「またこのセッションに参加したい」というように，音楽について相手と語り合う場に満足感を抱いていた．
実験後アンケートのPQ7（今後も相手と音楽の会話をしたい）も平均4.5と高スコアを記録し，曲を交互に聴き合い対話するというコミュニケーション形式が，相手との関係性を深める効果的な手段であることが示唆された．

% 「楽しかった」と答えた被験者は，音楽に対する思い入れや熱意がある人が多かった．曲に対しての自分のストーリーや物語がある人，聴いていた時の情景や時期を鮮明に覚えている場合などである．他にも，アーティスト自体への思い入れが強く，歌い方や音使いへのこだわりを話している人もいた．
一方で，「楽しめなかった」や「うまく話せなかった」と答えた被験者は，音楽への思い入れや過去の結びつきが少ないあるいは全くない人や，音楽について話す経験が不足している人であった．単純に曲が好きといった表面的な好みしか持たない人にとっては，深く話す内容がなく，会話を発展させることが難しかったという原因が考えられる．


\subsubsection{選曲時の意識}
選曲時の意識に関する実験結果から，重要な示唆が得られた．20人の被験者のうち14人は，選曲の基準として「思い入れや経験がある曲」「ストーリーとして話せそうな曲」など，個人的なエピソードと結びついた曲を重視していた．これは歌唱へのプレッシャーがない環境下で，より個人的な体験と結びついた曲を選びやすかったためと考えられる．
また，半数の被験者が「相手が知っていそうか」「好きそうな曲のジャンルか」という相手への配慮を示した．これは実際に相手に対して音楽を流すという形式が意識されていたことを示唆している．一方で，歌いやすさや普段のカラオケで歌う曲を意識した被験者は4人に留まった．

この結果は，本実験の設計に関する課題を明らかにしている．歌唱を伴わない音楽共有という形式は，通常のカラオケとは異なる選曲基準であった．実際のカラオケ環境では，曲を実際に歌う必要があり，歌唱力が試されるため自己開示の深さに影響を与える重要な要素となり得る．被験者からも「実際に歌わないから気楽に流すことができた」という意見が得られ，歌唱に対する心理的負担の有無が選曲基準に大きく影響することが確認された．

このことから，曲の自己開示レベルを構成する要素として設定した「歌いやすさ」は，本実験環境においては適切な指標ではなかった可能性が高い．今後は，歌唱を伴う実験との比較検証による選曲基準の違いの定量的分析が必要である．



\subsection{推薦システムについての考察}

本項では，推薦システムの提示方法や，段階的なアルゴリズムについての考察を行う．

\subsubsection{流れに合わせた選曲}
被験者からは「前の曲との流れを意識して選曲していた」「相手と同じアーティストの曲を入れた」といった意見があり，アーティストや楽曲間の関連性や，現在の雰囲気を考慮した推薦が自然な会話の展開を支援する可能性が示された．そのため，推薦システムの設計において，会話の流れと被験者の思い入れに応じた柔軟な選曲を支援することの重要性が明らかになった．
特に，推薦曲を強制的に選択させるのではなく，複数の選択肢から選べる提示方式の有効性が確認された．

ヒアリングの結果，最終フェーズでは「来て欲しい曲がこなかった」「今の雰囲気に合わなかった」という意見が得られた．これは，それまでの選曲の流れとは別に，個人的に共有したい楽曲が存在することを示している．
これらの知見から，自己開示を促進する音楽推薦システムには，現在の空気（会話の文脈や相手の選曲）に沿った推薦と個人の意図を反映できる柔軟性の両立が求められることが明らかになった．


\subsubsection{受容性を高める段階的なアプローチ}
HQ3（フェーズが進むにつれての楽曲選択意識の変化）への回答より，多くの被験者が前半で相手の知っている曲を意識的に選択していた．初期段階における共通点の活用が，音楽共有の受容性向上に効果的であることが明らかになった．これは，共通の楽曲が初期の会話の起点となり，「アニメ/ドラマの主題歌」といった共通の話題から会話を発展させやすいためと考えられる．
注目すべきは，4名の被験者から「共通の思い出のある曲から話が特に盛り上がった」という報告があった点である．「相手と初めて映画を見に行ったときの主題歌」（被験者S, T）や「2人でフェスに行ったときに感動した曲」（被験者I, J）といった，2人の思い出と結びついた楽曲が，受容性を大きく向上させることが示唆された．

これらの知見は，自己開示を促進する音楽共有において，共通知識や共有体験を基盤とした段階的なアプローチの有効性を示している．システム設計においては，初期段階で2人にとって関連のある楽曲を優先的に推薦することで，より自然な会話の展開が期待できる．